\chapter{Implementation and Tool Integration}
\label{sec:implementation}

This chapter should translate the high-level workflow into concrete implementation choices. It is the right place to explain which repository components implement the system, how tools are wired into the runtime, how artifacts are stored, and how analysis targets move through the harness.

\section{Repository Structure and Execution Path}

The implementation is split across the runtime package and the testing harness. The core workflow lives under \texttt{multi\_agent\_wf/}, while reproducible evaluation infrastructure lives under \texttt{Testing/}. This separation is itself important to the thesis because it distinguishes the research system from the measurement apparatus used to study it.

\begin{itemize}
\item \texttt{multi\_agent\_wf/config.py}: workflow configuration loading and normalization.
\item \texttt{multi\_agent\_wf/runtime.py}: tool loading, worker topology assembly, and run-time orchestration.
\item \texttt{multi\_agent\_wf/pipeline.py}: staged execution logic and stage contracts.
\item \texttt{multi\_agent\_wf/shared\_state.py}: intermediate artifact tracking and run-state mutation.
\item \texttt{Testing/harness/}: bundles, judging, sweeps, results aggregation, recovery, and viewers.
\end{itemize}

\begin{itemize}
\item \TODO{Describe why the runtime and evaluation harness are separated and what this buys for reproducibility.}
\item \TODO{Add a diagram that maps code modules to conceptual system responsibilities.}
\end{itemize}

\section{Workflow Configuration}

The workflow uses JSON-backed configuration for stage contracts, worker archetypes, architectures, prompts, and pipeline selection. That design is deliberate. It makes ablations and reproducible experiments easier.

\begin{itemize}
\item Pipeline presets determine which stages run and in what order.
\item Architecture presets determine the available worker mix.
\item Worker prompt shape, role-prompt mode, and review level tune behavior without rewriting the core runtime.
\end{itemize}

\begin{itemize}
\item \TODO{Explain how configuration is validated and loaded into Python structures.}
\item \TODO{Discuss why JSON workflow configuration is preferable to hard-coded stage logic for this research setting.}
\end{itemize}

\section{Ghidra Integration}

Ghidra is the primary structural-analysis backend in the system. The thesis should describe not only that it is used, but how it is used: prepared bundles, exported analysis artifacts, and agent-facing access through tool mediation rather than ad hoc manual interaction.

\begin{itemize}
\item Ghidra-backed analysis is prepared through bundle generation.
\item Agent stages consume Ghidra-derived artifacts. Where allowed, they can also query Ghidra-facing tools.
\item The system explicitly tracks whether analysis is occurring on the original sample or on a derived artifact such as an unpacked UPX binary.
\end{itemize}

\begin{itemize}
\item \TODO{Document the Ghidra bridge architecture and what metadata is exported into bundles.}
\item \TODO{Explain how pending Ghidra changes or mutation requests are controlled.}
\item \TODO{State which Ghidra capabilities are exposed to agents and which remain restricted.}
\end{itemize}

\section{Supporting Static and Dynamic Tooling}

The broader tool ecosystem matters because much of the thesis is about tool augmentation rather than raw model capability. The implementation chapter should make the tool landscape explicit.

\begin{itemize}
\item Static support tools include strings, FLOSS, and hash lookups.
\item Additional enrichers include capability extraction, YARA, binwalk, and UPX-oriented support.
\item Dynamic support infrastructure includes VM-oriented tooling and execution-time instrumentation where the task demands it.
\item Tool profiles are used to measure how much accuracy depends on the breadth of the tool ecosystem.
\end{itemize}

\begin{itemize}
\item \TODO{Describe how tool servers are grouped into static and dynamic partitions.}
\item \TODO{Explain how tool failures are surfaced back into run state and later analysis.}
\item \TODO{Clarify which tasks in the benchmark corpus actually require dynamic evidence.}
\end{itemize}

\section{Artifacts, Bundles, and Provenance}

Prepared bundles and run artifacts are central to the system's reproducibility story. Bundles capture analysis inputs and precomputed artifacts for a sample, while run records capture the outputs of one particular workflow configuration. The thesis should describe both layers.

\begin{itemize}
\item Bundles are reused when their generation inputs remain fresh.
\item Prompt or persona changes do not by themselves invalidate bundles.
\item Run records preserve final reports, judge outputs, live status, and per-task metadata.
\item Recovery and repaired-experiment workflows preserve provenance instead of mutating original results in place.
\end{itemize}

\begin{itemize}
\item \TODO{Explain the bundle lifecycle and why it is important for experiment cost and reproducibility.}
\item \TODO{Describe how recovery runs and rebuilt experiments preserve provenance.}
\item \TODO{Add one worked example of the files produced for a single sample-task run.}
\end{itemize}

\section{User-Facing Monitoring and Result Inspection}

The project also includes monitoring and browsing interfaces that expose live state and archived results. These interfaces are worth mentioning because they operationalize the system's provenance model and make failure analysis tractable.

\begin{itemize}
\item Live monitor for in-progress experiments.
\item Results browser for archived experiments, per-executable views, and detail panes.
\item Recovery tooling for failed-task reruns and rebuilt repaired experiments.
\end{itemize}

\begin{itemize}
\item \TODO{Briefly justify why these interfaces matter to the thesis instead of treating them as incidental engineering.}
\item \TODO{Show how run, task, and comparison state become inspectable by the experimenter.}
\end{itemize}
